\color{unidarkgreen}{\subsection{Логика отключения присоединений}}
\color{black}

\begin{enumerate}[label=\arabic{section}.\arabic{subsection}.\arabic{subsubsection}.\arabic{enumi}, labelsep=4pt, leftmargin=0pt, itemindent=57pt, itemsep=0pt, parsep=5pt] % если что можно крутить itemsep

\item
Ввод/вывод функции осуществляется либо посредством программного ключа \textbf{«ЛО Вn»}, либо с
использованием внешнего сигнала \textit{«Вывод действия на Вn»}.

%%%%%%%%%%%%%%%%%%%%%%%%%%%%%%%%%%%%%%%%%%%%%%%%%% Логика запрета АПВ %%%%%%%%%%%%%%%%%%%%%%%%%%%%%%%%%%%%%%%%%%%%%%%%%%
\subsubsection{Логика запрета АПВ}

\begin{enumerate}[label=\arabic{section}.\arabic{subsection}.\arabic{subsubsection}.\arabic{enumi}, labelsep=4pt, leftmargin=0pt, itemindent=57pt, itemsep=0pt, parsep=5pt] % если что можно крутить itemsep

\item
Сигналы запрета АПВ формируется в следующих условиях:
\begin{itemize}
\item от сигнала об аварийном снижении давления элегаза в ТТ;
\item при отключении при опробовании;
\item от общего сигнала запрета АПВ;
\item от команды оперативного отключения выключателя;
\item при срабатывании ЗНФ (при введенной уставку \textbf{«Запрет АПВ Вn от ЗНФ Вn»});
\item если выключатель не включился за время \textbf{«Tожид. усл. пуска АПВ Вn»} (при введенном программном ключе \textbf{«Огр. ожид. пуска АПВ Вn»});
\item если смежный выключатель (ведущий) включился на КЗ.
\end{itemize}

\medskip

\begin{figure}[H]
\centering
\includegraphics[width=0.8\textwidth,height=0.8\textheight,keepaspectratio]{\fbpath/761.9901 Логика отключения присоединений/img/LZ_APV.pdf}
\captionof{figure}{Функциональный блок логики запрета АПВ В1}\label{pict_lo:LZ_APV}
\end{figure}

\medskip

\end{enumerate}

%%%%%%%%%%%%%%%%%%%%%%%%%%%%%%%%%%%%%%%%%%%%%%%%%% Логика отключения %%%%%%%%%%%%%%%%%%%%%%%%%%%%%%%%%%%%%%%%%%%%%%%%%%
\subsubsection{Логика отключения}

\begin{enumerate}[label=\arabic{section}.\arabic{subsection}.\arabic{subsubsection}.\arabic*, labelsep=4pt, leftmargin=0em, itemindent=65pt, parsep=0pt]

\item
Функциональный блок логики отключения предназначен для подачи сигнала на отключение силового выключателя.

\item
Команда на отключение выключателя \textit{«Отключить Вn»} формируется:
\begin{itemize}
\item при отключении ошиновки (\textit{«Отключение СШ»});
\item при отключении в результате опробования присоединения;
\item от команды оперативного отключения выключателя (\textit{«РКО Вn»});
\item при действии УРОВ на отключение «своего» выключателя (\textit{«УРОВ Вn на себя»}).
\end{itemize}

\item
При необходимости может быть введен подхват команды на отключение от токовых реле,
установленных в цепях ЭМО, в этом случае сигнал на отключение удерживается в сработанном состоянии до
исчезновения тока через ЭМО (пока блок-контакты выключателя не разорвут цепи отключения). Ввод
осуществляется программным ключом \textbf{«Подхват откл. Вn от РТ ЭМО»}. 

\item
При наличии сигнала \textit{«Блок. отключения Вn»}, формируемого функциональным блоком контроля
силового выключателя отключение выключателя блокируется.

\item
Выдержка времени на возврат \textbf{«Тпрод откл. Вn»} задает время продления сигнала отключения.

\medskip

\begin{figure}[H]
\centering
\includegraphics[width=0.8\textwidth,height=0.8\textheight,keepaspectratio]{\fbpath/761.9901 Логика отключения присоединений/img/LO_Q1.pdf}
\captionof{figure}{Функциональный блок логики отключения В1}\label{pict_lo:LO_Q1}
\end{figure}

\medskip

\small
\begin{longtable}{|>{\centering\arraybackslash}m{5.3cm}|>{\centering\arraybackslash}m{3.3cm}|>{\centering\arraybackslash}m{4.2cm}|>{\centering\arraybackslash}m{1.8cm}|>{\centering\arraybackslash}m{1cm}|}
% Первая страница: полный заголовок
\caption{Перечень уставок логики отключения присоединений\hfill\vspace{-0.5\baselineskip}}\label{lo:tbl1}\\
\hline
\rowcolor{gray!30}
Параметр & Условное обозначение на схеме & Значение/ Диапазон & Единица измерения & Шаг \\
\hline
\endfirsthead

% Промежуточные страницы: «Продолжение таблицы»
\caption*{\hspace{3pt}\emph{Продолжение таблицы \ref{lo:tbl1}}\hfill\ }\vspace{-0.5\baselineskip} \\
\hline
\rowcolor{gray!30}
Параметр & Условное обозначение на схеме & Значение/ Диапазон & Единица измерения & Шаг \\
\hline
\endhead

%===m>ДЗО|-
%===m
\end{longtable}
\normalsize	


\end{enumerate}

\end{enumerate}